\begingroup\color{darkorange}
\onecolumn
\section*{Syntax of Related Approaches}
\label{app:syntax}

This appendix compares the rule syntax of the approaches discussed in the paper and shows how each translates into the syntax used here.
Our rules have the form $a_1\flpor\dots\flpor a_n\flpif{\flpweight}B$ with $n\geq 1$, a weight $\flpweight$, a disjunctor for the head that may be chosen per rule, and a body $B$ built from constants, (default-negated) literals and monotone connectives.
Tables~\ref{tab:syntax-heads} and~\ref{tab:syntax-bodies} list, for each approach, the rule form and the connectives that are allowed in heads and bodies, and how negation, weights and truth values are treated.
Table~\ref{tab:syntax-locations} says where the syntax, the connectives and the semantics are defined, and Table~\ref{tab:syntax-translation} gives the translation into our syntax and its status.
Page numbers are those printed in the papers, except for the report of Lukasiewicz and Straccia~(2007), which has its own numbering, the chapter of Mushthofa et al.~(2020), which is numbered relative to the chapter, and the preprints of Mushthofa et al.~(2014), Pelov and Truszczy\'nski~(2004), Sen et al.~(2022) and Kettmann et al.~(2025), which have no printed numbers and are cited by PDF page.
We write $\flpand_t$ and $\flpor_t$ for a generic t-norm and t-conorm, $\flpand_l$ and $\flpor_l$ for the conjunction and disjunction of Łukasiewicz logic, and $\flpand_g$ and $\flpor_g$ for those of Gödel logic (min and max); several generic connectives of one paper are written $\flpand_{t_i}$ and $\flpor_{t_i}$.

    {\scriptsize
        \renewcommand{\arraystretch}{1.25}
        \setlength{\tabcolsep}{3pt}
        \begin{longtable}{@{}p{0.15\textwidth}p{0.40\textwidth}p{0.42\textwidth}@{}}
            \caption{Rule form and head of the rules.}\label{tab:syntax-heads}                                                                                                                                                                                                                                                                                                                                                                    \\
            \toprule
            \textbf{Approach}                     & \textbf{Rule form}                                                                                                                                      & \textbf{Head and head disjunction}                                                                                                                                                                                                  \\
            \midrule\endfirsthead
            \toprule
            \textbf{Approach}                     & \textbf{Rule form}                                                                                                                                      & \textbf{Head and head disjunction}                                                                                                                                                                                                  \\
            \midrule\endhead
            \bottomrule\endfoot
            \cite{VanNieuwenborghDV06}            & $\alpha\leftarrow\beta$ with $|\alpha|\leq1$, $\beta$ a set of extended literals                                                                        & One extended literal, even a default-negated one; empty head is a constraint. No disjunction (future work).                                                                                                                         \\
            \cite{VanNieuwenborghDV07}            & $a\leftarrow\beta$, $a$ a literal or $\bot$                                                                                                             & One literal ($a$ or $\neg a$) or $\bot$ (constraint); no default negation in heads.                                                                                                                                                 \\
            \cite{JanssenSVD09}                   & $a\leftarrow\alpha$, residual implicator per rule, $\alpha$ any expression over acceptable functions                                                    & One proposition or a constant (constraint). No disjunction.                                                                                                                                                                         \\
            \cite{JanssenVSD12}                   & $a\leftarrow T(b_1,\dots,b_n)$, $T$ an arbitrary t-norm chosen per rule                                                                                 & One atom or a constant in $[0,1]$ (constraint). No disjunction.                                                                                                                                                                     \\
            \cite{JanssenSVD12core}               & $a\leftarrow f(b_1,\dots,b_n)$, $f$ increasing, residual implicator of a t-norm                                                                         & One atom; a constant head is allowed in the extension for constraints. No disjunction.                                                                                                                                              \\
            \cite{SchockaertJVD09,SchockaertJV12} & FASP: $a\leftarrow\alpha$; fuzzy equilibrium logic: arbitrary formulas, rules $\alpha\to\beta$                                                          & FASP: one atom or constant. Equilibrium logic: arbitrary formulas, so disjunction (t-conorm $\flpor_t$ or max chosen for the theory), conjunctions and nested rules; $a\leftarrow b\lor c$ is equivalent to two rules only for max. \\
            \cite{BlondeelSVD14}                  & $g(a_1,\dots,a_n)\leftarrow f(b_1,\dots,b_m,\mathbf{not}\,c_1,\dots,\mathbf{not}\,c_k)$                                                                 & Regular FASP: any composition $g$ of $\flpand_l,\flpor_l,\flpor_g,\flpand_g$ over literals and constants. Strict disjunctive FASP: only $a_1\flpor_l\dots\flpor_l a_n$.                                                             \\
            \cite{MushthofaSD14}                  & $\alpha\leftarrow\beta$, head and body expressions                                                                                                      & Any nesting of $\flpand_l,\flpor_l,\flpor_g,\flpand_g$ over literals and constants; typical heads $a\flpor_l b$ and $a\,\max\,b$; constant head is a constraint.                                                                    \\
            \cite{MushthofaSD15}                  & as above                                                                                                                                                & Same grammar, but the paper restricts the head to $\flpor_l$-disjunctions of literals; $\flpor_l$ is allowed in bodies and heads.                                                                                                   \\
            \cite{MushthofaSD20}                  & as above                                                                                                                                                & As above; one literal in the head is a normal rule, otherwise disjunctive ($\flpor_l$ or max).                                                                                                                                      \\
            \cite{AlvianoP13}                     & $a\leftarrow b_1\flpand_l\dots\flpand_l b_m\flpand_l\mathbf{not}\,b_{m+1}\flpand_l\dots\flpand_l\mathbf{not}\,b_n$                                      & One fuzzy atom (an atom or a numeric constant). No disjunction.                                                                                                                                                                     \\
            \cite{AlvianoP15}                     & $\alpha\leftarrow\beta$, $\alpha$ a head expression                                                                                                     & A chain of atoms with one connective from $\flpand_l$, $\flpor_l$, $\flpor_g$, $\flpand_g$; $n=1$ is atomic. Heads with $\flpor_l$ or max are disjunctive.                                                                          \\
            \cite{AlvianoAP17}                    & $\alpha\leftarrow\beta$                                                                                                                                 & One atom (a variable or a constant). No disjunction.                                                                                                                                                                                \\
            \cite{Lukasiewicz01}                  & $(A_1\lor\dots\lor A_l\leftarrow B_1,\dots,B_m,\mathbf{not}\,C_1,\dots,\mathbf{not}\,C_n)[c,1]$                                                         & Disjunction of atoms ($l\geq0$, empty head is an integrity clause), evaluated as max under positively correlated interpretations.                                                                                                   \\
            \cite{LukasiewiczS07}                 & $a_1\lor_{\flpor_{t_{1}}}\dots\lor_{\flpor_{t_{l-1}}}a_l\leftarrow_{\flpand_{t_{0}}}b_1\land_{\flpand_{t_{1}}}\dots\;\mathbf{not}_{\ominus}\,b_m\geq v$ & Non-empty disjunction of atoms, each disjunction symbol with its own t-conorm $\flpor_{t_{i}}$, so the connective may differ per rule and per position.                                                                             \\
            \cite{LukasiewiczS09}                 & $a\leftarrow_{\flpand_{t_{0}}}b_1\land_{\flpand_{t_{1}}}\dots\mathbf{not}_\ominus b_m$ with threshold $v$                                               & One atom. Disjunction strategies only occur in the formulas of the probabilistic part, not in rules.                                                                                                                                \\
            \cite{Mateis00}                       & $h_1\lor\dots\lor h_n\leftarrow^{[x,y]}b_1,\dots,b_k$, $0<x\leq y\leq1$                                                                                 & Non-empty disjunction of atoms evaluated by the t-conorm $S^{(p)}$ dual to a Schweizer--Sklar t-norm; the parameter $p$ is chosen once for the whole program.                                                                       \\
            \cite{Guller02}                       & $(l_1\mathbin{\dot\lor}\dots\mathbin{\dot\lor}l_n\leftarrow k_1\land\dots\land k_m,\;c)$                                                                & Strong disjunction of literals ($n\geq0$) interpreted as the join of a complete Boolean lattice; classical negation allowed in head literals.                                                                                       \\
            \cite{Guller2004procedural}           & $(h_1\lor\dots\lor h_n\leftarrow b_1\,\&\dots\&\,b_m,\;c)$                                                                                              & Non-empty disjunction of atoms ($n>0$) interpreted as the supremum of the residuated lattice, one connective for the whole program.                                                                                                 \\
            \cite{Saad2009extended}               & $l:\mu\leftarrow l_1:\mu_1,\dots,\mathbf{not}\,l_m:\mu_m$                                                                                               & One annotated literal. No disjunction.                                                                                                                                                                                              \\
            \cite{Saad2010disjunctive}            & $l_1:\mu_1\;\mathbf{or}\dots\mathbf{or}\;l_k:\mu_k\leftarrow l_{k+1}:\mu_{k+1},\dots,\mathbf{not}\,l_n:\mu_n$                                           & Disjunction of annotated literals, read existentially: at least one disjunct reaches its own annotation; every disjunct has its own annotation.                                                                                     \\
            \cite{NievesOC07,NievesOC13}          & $\alpha:a_1\lor\dots\lor a_m\leftarrow a_{m+1},\dots,\mathbf{not}\,a_n$, $\alpha$ in a finite lattice                                                   & Classical disjunction ($m\geq0$, $m=0$ is a constraint), the weight applies to the clause as a whole; no t-conorm, no weights per disjunct.                                                                                         \\
            \cite{NicolasGSL06}                   & $(c\leftarrow a_1,\dots,a_n,\mathbf{not}\,b_1,\dots,\mathbf{not}\,b_m,\alpha)$                                                                          & One atom. No disjunction (future work).                                                                                                                                                                                             \\
            \cite{BautersSDV14tplp}               & $\lambda:l_0;\dots;l_k\leftarrow l_{k+1},\dots,\mathbf{not}\,l_n$                                                                                       & Strong disjunction ``;'' is the maximum of the necessities of the disjuncts; weak disjunction is the necessity of the clause $l_0\lor\dots\lor l_k$.                                                                                \\
            \cite{BautersSDV14ijar}               & as above                                                                                                                                                & Strong reading (maximum of necessities) and, in the second reading, classical disjunction inside each subprogram.                                                                                                                   \\
            \cite{StracciaOD09}                   & $g(B_1,\dots,B_k)\leftarrow f(A_1,\dots,A_n)$ with computable connectors                                                                                & Arbitrary connector $g$; main example is the disjunction $\max(A,B)$ (lattice join); $\bot$ is the empty head.                                                                                                                      \\
            \cite{LoyerS09}                       & $A\leftarrow\varphi$                                                                                                                                    & One atom; several rules with the same head are combined by the join in the completion.                                                                                                                                              \\
            \cite{MadridO12}                      & $p\leftarrow p_1\ast\dots\ast p_m\ast\neg p_{m+1}\ast\dots\ast\neg p_n\;;\vartheta$                                                                     & One atom. No disjunction.                                                                                                                                                                                                           \\
            \cite{MedinaOV01}                     & $(A\leftarrow_iB),\vartheta$                                                                                                                            & One atom. No disjunction in heads (disjunctors only in bodies).                                                                                                                                                                     \\
            \cite{CornejoLM18}                    & $\langle p\leftarrow_i@[p_1,\dots,p_m,\neg p_{m+1},\dots,\neg p_n];\vartheta\rangle$                                                                    & One atom. No disjunction.                                                                                                                                                                                                           \\
            \cite{CornejoLM20}                    & $\langle p\leftarrow_i@^e[p_1,\dots,p_m;p_{m+1},\dots,p_n];\vartheta\rangle$ and constraints with weight $\top$                                         & One atom, or a constant $c$ (constraint). No disjunction.                                                                                                                                                                           \\
            \cite{CornejoLM20math}                & Multi-adjoint normal programs and core FASP $a\leftarrow f(l_1,\dots,l_n)$                                                                              & One atom. No disjunction (future work).                                                                                                                                                                                             \\
            \cite{HeyninckAB22arxiv}              & $\bigvee_{i=1}^np_i\leftarrow\psi$                                                                                                                      & Non-empty classical disjunction of atoms, read inclusively (a model needs at least one true head atom); four-valued semantics.                                                                                                      \\
            \cite{PelovT04}                       & $A_1\lor\dots\lor A_k\leftarrow L_1\land\dots\land L_m$, with monotone aggregate atoms                                                                  & Classical disjunction of atoms ($k\geq1$), handled by non-deterministic operators.                                                                                                                                                  \\
            \cite{RiegelEtAl20}                   & no rules: a graph of weighted formulas                                                                                                                  & Not applicable; $\flpor_t$ is an ordinary connective of weighted formulas.                                                                                                                                                          \\
            \cite{SenCRG22}                       & first-order $h\leftarrow b_1\land\dots\land b_m$ learned from templates                                                                                 & One predicate; disjunction only inside templates (LNN-$\lor$ nodes).                                                                                                                                                                \\
            \cite{Antic24}                        & $a_0\leftarrow^{w}a_1,\dots,a_k$ with weights $w$                                                                                                       & One neuron. No disjunction.                                                                                                                                                                                                         \\
            \cite{KettmannHS25}                   & $p\leftarrow^{\vartheta}_iB$                                                                                                                            & One atom (disjunctive heads are not covered).                                                                                                                                                                                       \\
            \cite{KettmannSHS26arxiv}             & $\ell\leftarrow B$ (weight-free)                                                                                                                        & One literal (atom or strongly negated atom); the disjunctive case is future work.                                                                                                                                                   \\
        \end{longtable}

        \begin{longtable}{@{}p{0.15\textwidth}p{0.31\textwidth}p{0.18\textwidth}p{0.17\textwidth}p{0.14\textwidth}@{}}
            \caption{Bodies, negation, weights and truth values.}\label{tab:syntax-bodies}                                                                                                                                                                                                                                                                                                                                                                                       \\
            \toprule
            \textbf{Approach}                                & \textbf{Body connectives}                                                                                                                          & \textbf{Negation}                                                                                  & \textbf{Weights}                                                                                & \textbf{Truth values}                                 \\
            \midrule\endfirsthead
            \toprule
            \textbf{Approach}                                & \textbf{Body connectives}                                                                                                                          & \textbf{Negation}                                                                                  & \textbf{Weights}                                                                                & \textbf{Truth values}                                 \\
            \midrule\endhead
            \bottomrule\endfoot
            \cite{VanNieuwenborghDV06,VanNieuwenborghDV07}   & Extended literals combined by one user-chosen t-norm, fixed for the whole framework; no t-conorms, max, aggregators or constants.                  & Default negation with any negator; classical $\neg$ as a literal with a consistency condition.     & None; rule degrees from an implicator, combined by an aggregator ($y$-models).                  & Complete lattice, propositional.                      \\
            \cite{JanssenSVD09}                              & Any expression over ``acceptable'' functions (monotone or antitone in each argument), e.g.\ t-norms, t-conorms, max, averages.                     & Any decreasing function; negative occurrences are replaced in the reduct.                          & Aggregator over rule satisfaction ($k$-answer sets).                                            & Bounded complete lattice.                             \\
            \cite{JanssenVSD12}                              & One arbitrary t-norm per rule over atoms, constants and $N(a)$; no t-conorms (explained in Sect.~7).                                               & Arbitrary negator $N$ on atoms.                                                                    & None; model needs rule degree $1$; weights as $a\leftarrow T(b,\alpha)$.                        & $[0,1]$.                                              \\
            \cite{JanssenSVD12core}                          & Any increasing function $f$ per rule, so t-norms, t-conorms, max, min; decreasing functions are simulated.                                         & $\flpnot  a$ with any negator; involutive one for the simulations.                                 & Aggregator over rule degrees in the extension.                                                  & Complete lattice.                                     \\
            \cite{SchockaertJVD09,SchockaertJV12}            & FASP: any function monotone or antitone in each argument; equilibrium logic: any combination of $\flpand_t,\flpor_t,\to$, strong negation and not. & Strong negation and default negation on arbitrary formulas, evaluated in two worlds.               & $\rho$-models for partial rule satisfaction.                                                    & Intervals in $[0,1]$, rational constants.             \\
            \cite{BlondeelSVD14}                             & Any composition of $\flpand_l,\flpor_l,\flpor_g,\flpand_g$; strict programs: $\flpand_l$-conjunctions of literals, naf-literals and constants.     & $\mathbf{not}\,l=1-l$ on literals; classical $\neg a$ as fresh literal with $I(l)+I(\neg l)\leq1$. & None; rule value $\min(1-\beta+\alpha,1)$.                                                      & $[0,1]\cap\mathbb{Q}$.                                \\
            \cite{MushthofaSD14,MushthofaSD15,MushthofaSD20} & Any nesting of $\flpand_l,\flpor_l,\flpor_g,\flpand_g$ over literals and constants, so $\flpor_l$ in bodies; fixed Łukasiewicz connectives.        & $\mathbf{not}$ on literals only, $1-x$; classical negation as separate literal.                    & None.                                                                                           & $[0,1]$ or $\{\frac{i}{k}\}$.                         \\
            \cite{AlvianoP13}                                & Conjunction $\flpand_l$ of atoms and naf-atoms; one fixed t-norm (Gödel, Łukasiewicz or product).                                                  & $\mathbf{not}\,b=1-b$ on atoms.                                                                    & None.                                                                                           & $[0,1]$.                                              \\
            \cite{AlvianoP15,AlvianoAP17}                    & Arbitrary nesting of $\flpand_l,\flpor_l,\flpor_g,\flpand_g$ and $\flpnot $, fixed connectives.                                                    & $\flpnot \alpha=1-\alpha$ on arbitrary expressions.                                                & None.                                                                                           & $[0,1]$ (AP17: also $L_n$).                           \\
            \cite{Lukasiewicz01}                             & Conjunction of atoms and naf-atoms, evaluated as min; no choice of t-norm.                                                                         & Default negation $\mathbf{not}\,C$, read classically in the clause.                                & Lower probability bound $c$ per rule; model iff $\max(\text{head},C\text{'s})\geq c-1+\min(B)$. & Finite chain $\{0,\frac1{n-1},\dots,1\}$.             \\
            \cite{LukasiewiczS07,LukasiewiczS09}             & Conjunction of atoms, constants and naf-atoms; each $\land$ with its own t-norm, per rule and position; no disjunction in bodies.                  & $\mathbf{not}_\ominus$ with its own negation strategy on atoms.                                    & Threshold $v$, conjoined to the body with $\flpand_{t_{0}}$.                                    & $[0,1]$, function-free first order.                   \\
            \cite{Mateis00}                                  & Conjunction of atoms by the t-norm $T^{(p)}$ (global); no negation.                                                                                & None.                                                                                              & Interval $[x,y]$ combined by the product t-norm.                                                & Intervals in $[0,1]$.                                 \\
            \cite{Guller02,Guller2004procedural}             & Conjunction (meet in 2002, t-norm $\ast$ in 2004); no disjunction.                                                                                 & 2002: classical negation only; 2004: none.                                                         & Degree $c$ per implication; $(\varphi,c)$ true iff $\|\varphi\|\geq c$.                         & Complete Boolean (2002) or residuated (2004) lattice. \\
            \cite{Saad2009extended,Saad2010disjunctive}      & Conjunction of annotated literals, satisfied iff every item reaches its annotation; fixed max/min.                                                 & Classical $\neg$ ($1-x$) and default $\mathbf{not}$, with partial interpretations.                 & Annotations on every literal; no rule weight.                                                   & $[0,1]$.                                              \\
            \cite{NievesOC07,NievesOC13}                     & Conjunction of atoms and naf-literals, classical two-valued; weights combine by glb in the proof system.                                           & Boolean naf via reduct; strong negation by fresh atoms.                                            & Necessity degree per clause in a finite (possibly partial) lattice.                             & Two-valued atoms.                                     \\
            \cite{NicolasGSL06}                              & Conjunction by min of necessity degrees.                                                                                                           & Boolean naf via a Gödel-style reduct.                                                              & Necessity $\alpha$ per rule.                                                                    & Finite chain $N\subseteq\,]0,1]$.                     \\
            \cite{BautersSDV14tplp,BautersSDV14ijar}         & Conjunction by min of necessities; weak case: clauses allowed in bodies.                                                                           & Classical $\neg$ and naf, with Gödel ($\mathrm{PASP}_G$) or Łukasiewicz ($1-V$) negator.           & Certainty $\lambda$ per rule as a constraint on possibility distributions.                      & Valuations in $[0,1]$ over literals.                  \\
            \cite{StracciaOD09}                              & Any computable connector, also non-monotone (e.g.\ $1-x$, $\max(0,x+y-1)$, weighted sums).                                                         & No default negation, no stable semantics.                                                          & None.                                                                                           & Arbitrary complete lattice.                           \\
            \cite{LoyerS09}                                  & Any continuous function in a family containing meet and join.                                                                                      & Involutive order-reversing $\neg$.                                                                 & None.                                                                                           & Complete lattice; intervals $L\times L$.              \\
            \cite{MadridO12}                                 & $\ast$-conjunction of atoms and naf-atoms, one fixed $\ast$; no disjunction.                                                                       & Any decreasing $\neg$ with $\neg0=1$, $\neg1=0$.                                                   & Weight $\vartheta$, $I(B)\ast\vartheta\leq I(p)$.                                               & Residuated lattice.                                   \\
            \cite{MedinaOV01}                                & Conjunctors, disjunctors and aggregators, several adjoint pairs.                                                                                   & None.                                                                                              & Confidence $\vartheta$, adjoint pair per rule.                                                  & Bounded complete upper semilattice.                   \\
            \cite{CornejoLM18,CornejoLM20,CornejoLM20math}   & One aggregator per rule; extended aggregators monotone in the first and antitone in the other arguments; core FASP: any monotone $f$.              & Negation operator on atoms; several negations via antitone arguments (EMALP).                      & Weight $\vartheta$ with adjoint pairs (not in core FASP).                                       & Complete lattice.                                     \\
            \cite{HeyninckAB22arxiv}                         & General propositional formulas over $\neg,\land,\lor$ (Belnap); conjunctions of literals for the stable semantics.                                 & Default negation $\neg$ with an involution.                                                        & None.                                                                                           & Belnap's four-valued bilattice.                       \\
            \cite{PelovT04}                                  & Conjunction of literals; literals are atoms or monotone aggregate atoms.                                                                           & Default negation on atoms and aggregates.                                                          & None.                                                                                           & Boolean, four-valued pairs.                           \\
            \cite{RiegelEtAl20}                              & Weighted Łukasiewicz connectives $\flpand_l^{w}$, $\flpor_l^{w}$, $\to$, Gödel and product variants; monotone activation functions.                & $\neg x=1-x$.                                                                                      & Weights $w_i\geq0$ and bias $\beta\geq0$ per connective.                                        & $[0,1]$ with bounds $[L,U]$.                          \\
            \cite{SenCRG22}                                  & $\mathrm{LNN}\text{-}\land(x;\beta,w)=\mathrm{relu}_1(\beta-w^{\top}(1-x))$ and its dual.                                                          & $\neg x=1-x$.                                                                                      & Learnable $\beta$, $w\geq0$.                                                                    & $[0,1]$ with threshold $\alpha$.                      \\
            \cite{Antic24}                                   & Weighted threshold units $\sum w\cdot I(b)\geq\theta$; no connectives.                                                                             & Negative weights act like negation.                                                                & Real weights and thresholds.                                                                    & Boolean.                                              \\
            \cite{KettmannHS25,KettmannSHS26arxiv}           & Monotone connectives (conjunctors, disjunctors, aggregators) and constants.                                                                        & $\flpnot $ antitone involution (2025: on atoms; 2026: also classical $\neg$).                      & Weight $\vartheta$ with adjoint pairs (2025); weight-free (2026).                               & Complete lattice; paraconsistent pairs in 2026.       \\
        \end{longtable}

        \begin{longtable}{@{}p{0.18\textwidth}p{0.58\textwidth}p{0.20\textwidth}@{}}
            \caption{Translation into our syntax $a_1\flpor\dots\flpor a_n\flpif{\flpweight}B$.}\label{tab:syntax-translation}                                                                                                                                                                                                                                                                                             \\
            \toprule
            \textbf{Approach}                                                                               & \textbf{Translation}                                                                                                                                                                                                                                                & \textbf{Status}                        \\
            \midrule\endfirsthead
            \toprule
            \textbf{Approach}                                                                               & \textbf{Translation}                                                                                                                                                                                                                                                & \textbf{Status}                        \\
            \midrule\endhead
            \bottomrule\endfoot
            \cite{Guller2004procedural,Guller02}                                                            & Head disjunction as lattice join, weight $\flpweight=c$, body conjunction by the t-norm $\ast$ (2002: literals and the lattice meet). Graded disjunctions are consequences, not rules.                                                                              & exact                                  \\
            \cite{Saad2010disjunctive}                                                                      & Body literal $l:\mu$ becomes $\gecon{\mu}(l)$ and $\mathbf{not}\,l:\mu$ becomes $\flpnot\gecon{\mu}(l)$; head $l_1\flpor\dots\flpor l_k\flpif{\mu}\dots$ if all head annotations equal $\mu$. Different annotations per disjunct need an existential head operator. & partial                                \\
            \cite{NievesOC07,NievesOC13}                                                                    & $a_1\flpor\dots\flpor a_n\flpif{\alpha}\bigwedge B^+\land\bigwedge\flpnot b$ with join and meet of $Q$ and a Gödel negator. Their two layers (an answer set of the projection first, then weights) are not reproduced.                                              & lossy                                  \\
            \cite{BautersSDV14tplp}                                                                         & Strong disjunction: weight $\flpweight=\lambda$, head disjunction max, body min, $\flpnot x=1-x$. Weak (clausal) disjunction is not expressible.                                                                                                                    & exact for the rule form, none for weak \\
            \cite{BautersSDV14ijar}                                                                         & As the strong case; the weights as uncertainty about the rules (possibility over subprograms) are not truth-functional.                                                                                                                                             & none for the second reading            \\
            \cite{Lukasiewicz01}                                                                            & Head $A_1\flpor\dots\flpor A_l$ with max, body min, weight $c$ via $c\flpand_l x=\max(c+x-1,0)$. Constants in the head of the reduct become $t\ominus d$ in the body. Model states have no counterpart in the syntax.                                               & partial                                \\
            \cite{LukasiewiczS07}                                                                           & Weight $v$, adjoint conjunctor $\flpand_{t_{0}}$, head disjunctor $\flpor_t$ per rule, body conjunction $\flpand_{t_{i}}$, $\flpnot=\ominus$. DL atoms are not part of our language.                                                                                & exact for rules                        \\
            \cite{Mateis00}                                                                                 & Head with the t-conorm $S^{(p)}$, body with the t-norm $T^{(p)}$, weights combined by the product; interval weights $[x,y]$ and interval-valued interpretations need a lower and an upper program ($\flpweight=x$, $\flpweight=y$), which is a different semantics. & lossy                                  \\
            \cite{MushthofaSD14,MushthofaSD15,MushthofaSD20,BlondeelSVD14}                                  & Head $a_1\lkor\dots\lkor a_n$ with weight $1$; facts $a\leftarrow\lkc{c}$ become $a\flpif{c}\top$; max heads use the join; constraints use the encoding of Kettmann et al.\ with fresh atoms. Heads with $\flpand_l$ are outside our syntax.                        & exact for $\flpor_l$ and max heads     \\
            \cite{AlvianoP15}                                                                               & As above; $\flpand_l$-heads and max-heads are shifted by auxiliary Booleans (their eqs.~(5)--(7)).                                                                                                                                                                  & exact for $\flpor_l$ and max           \\
            \cite{SchockaertJVD09,SchockaertJV12}                                                           & Heads that are disjunctions of atoms translate directly; guess atoms $a\flpor_l\neg a\leftarrow$; arbitrary head formulas and the (here, there) semantics have no counterpart.                                                                                      & partial                                \\
            \cite{StracciaOD09}                                                                             & Heads $g=\flpor_l$ or join of atoms give our heads; non-monotone body connectors and general head connectors are not covered.                                                                                                                                       & partial                                \\
            \cite{VanNieuwenborghDV06,VanNieuwenborghDV07}                                                  & $a\leftarrow_1 T(b_1,\dots,b_n)$ with naf as $\flpnot$ and classical negation as literals. The aggregated rule degrees ($y$-models) are not representable.                                                                                                          & lossy                                  \\
            \cite{JanssenSVD09,JanssenVSD12,JanssenSVD12core,AlvianoP13,AlvianoAP17,LoyerS09}               & $a\flpif{}B$ with the body as is; weights as $a\flpif{\flpweight}b$. Already covered by Kettmann et al.~(2025).                                                                                                                                                     & exact                                  \\
            \cite{NicolasGSL06,MadridO12,CornejoLM18,CornejoLM20,CornejoLM20math,MedinaOV01,LukasiewiczS09} & $a\flpif{\flpweight}B$ with the aggregator or weight kept as is; constraints through fresh atoms. Covered by Kettmann et al.~(2025, 2026).                                                                                                                          & exact                                  \\
            \cite{HeyninckAB22arxiv,PelovT04}                                                               & Crisp special case with $L=\{0,1\}$ and head disjunction $\vee$; aggregates are not covered.                                                                                                                                                                        & exact (Boolean)                        \\
            \cite{RiegelEtAl20,SenCRG22}                                                                    & Tight programs and the input-aware completion $a\leftrightarrow B_1\lor\dots\lor B_n$; weights and biases of the connectives are dropped.                                                                                                                           & partial                                \\
            \cite{Antic24}                                                                                  & $a\flpif{1}\gecon{\theta}(\ldots)$ with a threshold connective over the weighted sum.                                                                                                                                                                               & exact (Boolean)                        \\
            \cite{KettmannHS25,KettmannSHS26arxiv}                                                          & Our syntax with $n=1$.                                                                                                                                                                                                                                              & exact                                  \\
        \end{longtable}

        \begin{longtable}{@{}p{0.17\textwidth}p{0.20\textwidth}p{0.58\textwidth}@{}}
            \caption{Operators and semantic notions for disjunctive heads.}\label{tab:operators}                                                                                                                                                                                                                                                                                                                                                                                                                                                                                                                                                    \\
            \toprule
            \textbf{Approach}                                                                        & \textbf{Notion}                                          & \textbf{Definition}                                                                                                                                                                                                                                                                                                                                                                                                                                                               \\
            \midrule\endfirsthead
            \toprule
            \textbf{Approach}                                                                        & \textbf{Notion}                                          & \textbf{Definition}                                                                                                                                                                                                                                                                                                                                                                                                                                                               \\
            \midrule\endhead
            \bottomrule\endfoot
            \cite{BlondeelSVD14,MushthofaSD14,MushthofaSD15,MushthofaSD20,AlvianoP15,LukasiewiczS07} & Reduct, answer set                                       & $P^I$ replaces every $\mathbf{not}\,\alpha$ by the constant $1-I(\alpha)$ ($\ominus_iI(b_i)$ in LS07). $I$ is an answer set iff $I\models P$ and no $J<I$ satisfies $P^I$ (minimal model of $P^I$).                                                                                                                                                                                                                                                                               \\
            \cite{MushthofaSD14}                                                                     & Grid translation $\mathit{Tr}(P,k)$                      & Level atoms $a_i$ with $a_i\equiv I(a)\geq\frac ik$, rules $a_i\leftarrow a_{i+1}$; head $a\,\max\,b\leftarrow c$ becomes $a_i\lor b_i\leftarrow c_i$; head $a\flpor_l b\leftarrow c$ becomes $p_{0,i}\lor\dots\lor p_{i,0}\leftarrow c_i$ with $a_s\leftarrow p_{s,t}$, $b_t\leftarrow p_{s,t}$.                                                                                                                                                                                 \\
            \cite{MushthofaSD15,AlvianoP15}                                                          & Shift                                                    & $a_1\flpor_l\dots\flpor_l a_n\leftarrow\beta$ becomes $a_i\leftarrow\beta\flpand_l\bigwedge\nolimits_{t,\,j\neq i}\flpnot a_j$ for all $i$; faithful for HCF programs (AP15) and for SRCF programs (MSD15).                                                                                                                                                                                                                                                                       \\
            \cite{LukasiewiczS07}                                                                    & Crisp transform                                          & Over $\mathit{TV}_n$ (Zadeh only): atoms $a_\alpha$ for $\alpha\neq0$, rules $a^1_\alpha\lor\dots\lor a^l_\alpha\leftarrow b^1_\alpha,\dots,\mathbf{not}\,b^m_\gamma$ with $\gamma=1-\alpha+\frac1n$ for all $\alpha\leq v$, and $p_{\alpha-1/n}\leftarrow p_\alpha$.                                                                                                                                                                                                             \\
            \cite{Guller2004procedural}                                                              & $\mathcal{C}_P$, $\mathit{DS}$, $\mathit{FS}$            & $(D,c)\in\mathcal{C}_P(I)$ iff $D=D_0\cup\bigcup_i(D_i\setminus\{b_i\})$ and $c=\vartheta\ast c_1\ast\dots\ast c_m$ for a rule $D_0\leftarrow_\vartheta b_1,\dots,b_m$ and $(D_i,c_i)\in I$ with $b_i\in D_i$; $\mathit{FS}(P)$ closes $\mathrm{lfp}(\mathcal{C}_P)$ under supersets and lower degrees; $\mathit{DS}(P)=\mathit{FS}(P)$.                                                                                                                                          \\
            \cite{Lukasiewicz01}                                                                     & Model, $T^M_P$, $T^s_P$                                  & $p\models C$ iff $\max(\max_ip(A_i),\max_jp(C_j))\geq c-1+\min_kp(B_k)$. $T^M_P(\mathcal P)=\bigcup_{p\in\mathcal P}\mathrm{models}_p(\mathrm{state}_{\mathcal P}(p))$ on up-closed sets of interpretations (lfp at $\omega$, $\mathit{MM}(P)=\mathrm{can}(\mathrm{lfp})$). $T^s_P(S)=\exp\{A^\alpha_1\lor\dots\lor A^\alpha_l\lor D_1\lor\dots\lor D_m\mid B_i^{\beta_i}\lor D_i\in S,\ \alpha=c-1+\min\beta_i\}$. Stable: minimal model of $P/q$ (head constant $\max q(C_j)$). \\
            \cite{Mateis00}                                                                          & $p$-model, computation                                   & $I(H(r))\geq T_2(I(B(r)),[x,y])$ with $T_2$ the product; a computation step picks an unsatisfied rule $r$ and $A\in H(r)$ and sets $I(A):=T_2(I(B(r)),[x,y])$; for $p=+\infty$ the choice tree is finite.                                                                                                                                                                                                                                                                         \\
            \cite{StracciaOD09}                                                                      & $T_P$, $F_P$, orbit                                      & $T_P(I)=\{J\mid J\models P[I],\,I\leq J\}$, $F_P(I)=\{J\mid J\models P[I]\}$; an orbit starts at $\bot$, takes $x_{\alpha+1}\in T_P(x_\alpha)$ and suprema at limits; minimal fixed points are limits of orbits.                                                                                                                                                                                                                                                                  \\
            \cite{Saad2010disjunctive}                                                               & Models, reduct                                           & $I\models l:\mu$ iff $\mu\leq I(l)$; head satisfied iff some disjunct is; reduct drops rules with an unsatisfied $\mathbf{not}\,l:\mu$ and the negative literals; answer set = minimal model of the reduct.                                                                                                                                                                                                                                                                       \\
            \cite{NievesOC13}                                                                        & $P_M$, $\mathrm{GPPE}$, $T$, $\mathit{Sem}_{\min}$       & $P_M=\{n(r):A\cap M\leftarrow B^+\mid A\cap M\neq\emptyset,\,B^-\cap M=\emptyset,\,B^+\subseteq M\}$; $\mathrm{GPPE}(\alpha:A\leftarrow B^+\cup\{b\},\ \alpha_1:A_1)=\min(\alpha,\alpha_1):A\cup(A_1\setminus\{b\})\leftarrow B^+$ for $b\in A_1$; $T(P)=P\cup\{\mathrm{GPPE}(r_1,r_2)\}$; $\mathit{Sem}_{\min}(P)=\{(x,\bigvee\{\alpha\mid\alpha:x\in P\})\}$; answer set $M=\mathit{Sem}_{\min}(\Pi(P_{M^*}))$ for a classical answer set $M^*$.                                \\
            \cite{BautersSDV14tplp}                                                                  & Strong answer set                                        & Constraint $\max(N(l_0),\dots,N(l_k))\geq\min(N(l_{k+1}),\dots,1-V(l_{m+1}),\dots,\lambda)$; $V$ is an answer set iff some minimally specific possibility distribution $\pi$ has $N_\pi=V$.                                                                                                                                                                                                                                                                                       \\
            \cite{SchockaertJV12}                                                                    & Equilibrium model                                        & $V:\{h,t\}\times\mathit{At}\to$ subintervals of $[0,1]$ with $V(h,a)\supseteq V(t,a)$; $V(h,\mathbf{not}\,\alpha)=[1-V^-(t,\alpha),1-V^-(h,\alpha)]$; model iff $V^-(h,\alpha)=1$; $V\preceq V'$ iff $V(t,\cdot)=V'(t,\cdot)$ and $V(h,\cdot)\supseteq V'(h,\cdot)$; equilibrium model: $\preceq$-minimal with $V(h,\cdot)=V(t,\cdot)$.                                                                                                                                           \\
            \cite{HeyninckAB22arxiv}                                                                 & $\mathit{IC}_P$, stable, WF state                        & $\mathit{IC}^l_P(x,y)=\{x_1\subseteq\bigcup HD^l\mid\forall\Delta\in HD^l\colon x_1\cap\Delta\neq\emptyset\}$ with $HD^l$ the heads of rules whose body holds at $(x,y)$; $\mathit{IC}^u_P(x,y)=\mathit{IC}^l_P(y,x)$; stable operator $\mathrm{C}(\mathit{IC}^l_P)(y)=\min\{x\mid x\in\mathit{IC}^l_P(x,y)\}$; well-founded state $=\mathrm{KK}(\mathrm{S}(\mathit{IC}_P)')$.                                                                                                    \\
            \cite{PelovT04}                                                                          & $N_P^{\mathit{Sel}}$                                     & $N_P^{\mathit{Sel}}(I)=\mathit{Sel}(HR_P(I))$ with $HR_P(I)$ the heads of rules whose body holds in $I$; $\mathit{Mod}(P)=\mathrm{pre}(N_P)$; stable iff $I\in\mathrm{mfp}(A_P(\cdot,I))$.                                                                                                                                                                                                                                                                                        \\
            \cite{LoyerS02}                                                                          & $T_P$, $S_P$, $\mathit{CS}$ (normal programs with $f_d$) & $T_P(I,J)(A)=f_d(\{f_p(\alpha_r,f_c(I(B_i),\neg J(C_j)))\mid r\text{ has head }A\})$; $S_P(J)=\mathrm{lfp}\,T_P(\cdot,J)$; $S^\bot_P=(S_P\circ S_P)^\infty(I_\bot)$, $S^\top_P=(S_P\circ S_P)^\infty(I_\top)$; $\mathit{CS}(P)=S^\bot_P\cap S^\top_P$.                                                                                                                                                                                                                            \\
            Ours                                                                                     & $N$, stable, KK and WF state                             & $N_l(L,U)=\{J\mid\forall r\colon J(a_1)\mathbin{\flpcontf{\flpor_r}}\dots\mathbin{\flpcontf{\flpor_r}}J(a_n)\geq\flpweight_r\mathbin{\flpcontf{\flpand}}(L,U)^{B_r},\ J\leq\mathrm{cap}\}$, $N_u(L,U)=N_l(U,L)$; $(I,I)$ stable iff $I\in\min N_l(\cdot,I)$ iff $I$ is a minimal model of $P^I$; KK and WF states as in HAB.                                                                                                                                                      \\
        \end{longtable}
    }

\subsection*{Differences between the frameworks}

\paragraph{Head disjunction.}
The approaches read a head disjunction in different ways.
(i)~\emph{None}: normal programs only (all multi-adjoint, core FASP, Loyer--Straccia, Nicolas et al., Saad 2009, our earlier papers).
(ii)~\emph{Lattice join or max}: Guller, Lukasiewicz (2001), the strong disjunction of Bauters et al., and the classical disjunction of Nieves et al., Heyninck et al.\ and Pelov and Truszczy\'nski; on a chain a disjunction $a\lor b\geq x$ holds iff $a\geq x$ or $b\geq x$.
(iii)~\emph{Łukasiewicz t-conorm}: Blondeel et al.\ and Mushthofa et al.\ with a fixed $\flpor_l$; Alviano and Pe\~naloza (2015) allow $\flpand_l$, $\flpor_l$, $\flpor_g$ or $\flpand_g$ as a head connective (one per rule).
(iv)~\emph{t-conorm chosen per rule or per symbol}: only Lukasiewicz and Straccia (2007), with arbitrary t-conorms.
(v)~\emph{One family chosen once}: Mateis (parameter $p$ of the Schweizer--Sklar family) and Guller (the lattice join of the residuated lattice).
(vi)~\emph{Existential per-disjunct thresholds}: Saad, where each disjunct has its own annotation.
(vii)~\emph{Clausal}: the weak disjunction of Bauters et al., where the clause itself has a necessity.
(viii)~\emph{Arbitrary head connector or formula}: Straccia et al.\ and fuzzy equilibrium logic.
Only (ii) with max has finitely many minimal head satisfiers on a chain; with the Łukasiewicz or product t-conorm they form a continuum.

\paragraph{Body connectives.}
One fixed t-norm for the whole program is assumed by Van Nieuwenborgh et al., Alviano and Pe\~naloza (2013), Madrid and Ojeda-Aciego, Mateis and Guller.
One t-norm per rule is allowed by Janssen et al.\ (2012), and one per position by Lukasiewicz and Straccia.
Any increasing function is allowed in the core FASP languages and by Cornejo et al.\ (aggregators); Medina et al.\ allow conjunctors, disjunctors and aggregators with several adjoint pairs; Straccia et al.\ even allow non-monotone connectors.
Disjunction in bodies ($\flpor_l$ or max) is allowed by Blondeel et al., Mushthofa et al., Alviano and Pe\~naloza, the core FASP papers and the multi-adjoint papers, but not by Saad, Guller, Mateis, Nieves et al., Bauters et al.\ (strong) and Lukasiewicz and Straccia.
It matters semantically: recursive $\flpor_l$ in bodies breaks completion and loop formulas (Janssen et al.~2012) and the shift of a disjunctive head (Mushthofa et al.).

\paragraph{Negation.}
Default negation uses $1-x$ (Mushthofa et al., Blondeel et al., Alviano and Pe\~naloza, Bauters et al.\ for $\mathrm{PASP}_L$), a Gödel-type negator (Nicolas et al., Nieves et al., Bauters et al.\ for $\mathrm{PASP}_G$) or an arbitrary decreasing function (Janssen et al., Madrid and Ojeda-Aciego, Cornejo et al., Lukasiewicz and Straccia, Kettmann et al.).
Saad, Guller (2002), Bauters et al.\ and Nieves et al.\ also have classical negation; it is a literal in Saad and Bauters et al., a fresh atom in Nieves et al.\ and Mushthofa et al., and a pair component in Kettmann et al.\ (2026).
Guller (2004), Mateis, Medina et al.\ and Straccia et al.\ have no negation.

\paragraph{Weights.}
Rule weights with adjoint pairs: Medina et al., Madrid and Ojeda-Aciego, Cornejo et al., Kettmann et al., Guller and Lukasiewicz and Straccia.
Weights as necessity or certainty: Nicolas et al., Nieves et al., Bauters et al.; as probability bounds or intervals: Lukasiewicz (2001), Mateis.
Annotations on literals instead of rules: Saad.
No weights, rule degrees aggregated: Van Nieuwenborgh et al., Janssen et al.\ (aggregator expressions).
No weights at all: Mushthofa et al., Blondeel et al., Alviano and Pe\~naloza, Loyer and Straccia, Straccia et al.

\paragraph{Truth values and models.}
Most works use $[0,1]$ (some with rational constants or finite grids $L_k$); Guller, Medina et al., Cornejo et al.\ and Madrid and Ojeda-Aciego use (residuated) complete lattices; Nieves et al.\ use a finite, possibly partially ordered lattice of weights over two-valued atoms; Mateis uses intervals; Heyninck et al.\ use Belnap's four values.
The semantics is a minimal model of a reduct in most disjunctive fuzzy approaches (Lukasiewicz and Straccia, Blondeel et al., Mushthofa et al., Alviano and Pe\~naloza, Saad), a fixpoint over graded disjunctions in Guller, and a two-layer possibilistic semantics in Nieves et al.\ and Bauters et al.

\paragraph{Consequence for our syntax.}
A uniform syntax needs a head disjunctor chosen per rule (to cover (ii)--(v)), cannot express per-disjunct annotations (vi) or clausal necessity (vii) with one weight, and leaves out heads with $\flpand_l$ and arbitrary head formulas (viii).
Rules with a weight, body connectives from monotone families and default negation are those of Kettmann et al., so the normal-program approaches are covered by our earlier papers.

\subsection*{Running example}
Let $\lset=[0,1]$ with the Gödel connectives, i.e.\ head disjunction $\flpor_g$ (maximum) and conjunction and weight combination $\flpand_g$ (minimum), and let $P$ consist of $r_1:\ a\flpor b\flpif{0.8}\top$, $r_2:\ c\flpif{}a$ and $r_3:\ c\flpif{}b$ (weight $1$ omitted). Triples are $(a,b,c)$.

    {\scriptsize
        \renewcommand{\arraycolsep}{2pt}
        \begin{longtable}{@{}p{0.20\textwidth}p{0.50\textwidth}p{0.25\textwidth}@{}}
            \caption{Fixpoint constructions on the running example.}\label{tab:example}                                                                                                                                                                                                                                                                                                                                                                \\
            \toprule
            \textbf{Operator}                                                                                                                                                                    & \textbf{Construction}                                                                                                                                         & \textbf{Result}                                                                     \\
            \midrule\endfirsthead
            \toprule
            \textbf{Operator}                                                                                                                                                                    & \textbf{Construction}                                                                                                                                         & \textbf{Result}                                                                     \\
            \midrule\endhead
            \bottomrule\endfoot
            Models of $P$                                                                                                                                                                        & $\max(a,b)\geq0.8$, $c\geq a$, $c\geq b$                                                                                                                      & minimal: $(0.8,0,0.8)$, $(0,0.8,0.8)$                                               \\
            Answer sets \cite{BlondeelSVD14,MushthofaSD15,AlvianoP15,LukasiewiczS07}, equilibrium models \cite{SchockaertJV12}, Bauters strong \cite{BautersSDV14tplp}, Saad (equal annotations) & minimal models of $P^I=P$; equilibrium: $V(h,\cdot)=V(t,\cdot)=[I,1]$                                                                                         & the two minimal models; Bauters: $\{a^{.8},c^{.8}\}$, $\{b^{.8},c^{.8}\}$           \\
            Ours, $N_l$                                                                                                                                                                          & $N_l(\bot)$: $\min=\{(.8,0,0),(0,.8,0)\}$; then $\{(.8,0,.8),(0,.8,.8)\}$; fixed. Upper: $\downarrow\!\{(.8,.8,1)\}$, then $\downarrow\!\{(.8,.8,.8)\}$       & KK state $(\uparrow\!\{(.8,0,.8),(0,.8,.8)\},\ \downarrow\!\{(.8,.8,.8)\})$         \\
            Ours, stable and WF                                                                                                                                                                  & $C(N_l)(y)=\min\mathrm{mod}(P)$ for all $y$; one step from the KK state                                                                                       & stable $(I,I)$ for both minimal models; WF $=(\min\!\uparrow,\ \min\!\downarrow)$   \\
            Shifted approximator (draft)                                                                                                                                                         & lower bound of $a$: $\bigglb\{z\mid z\flpcontf{\flpor_g}U(b)\geq0.8\}=0$ because $U(b)=1$; upper bound of $a,b$: $0.8$; of $c$: $U(a)\flpcontf{\flpor_g}U(b)$ & KK $=((0,0,0),(.8,.8,.8))$: $c$ undecided; stable fixpoints: the two minimal models \\
            Straccia et al. $T_P$ \cite{StracciaOD09}                                                                                                                                            & $x_0=\bot$; $T_P(x_0)\ni(.8,0,0),(0,.8,0)$; then $(.8,0,.8)$ resp.\ $(0,.8,.8)$                                                                               & two orbits, two minimal models                                                      \\
            Mateis, $p=+\infty$ \cite{Mateis00}                                                                                                                                                  & root $\bot$; raise $a$ or $b$ to $0.8$; then raise $c$ to $0.8$                                                                                               & choice tree of depth $2$, same leaves                                               \\
            Lukasiewicz $T^M_P$ \cite{Lukasiewicz01}                                                                                                                                             & grid $\{0,.2,\dots,1\}$; minimal elements of $T^M\!\uparrow k$: $\{(.8,0,0),(0,.8,0)\}$, then $\{(.8,0,.8),(0,.8,.8)\}$                                       & $\mathit{MM}(P)$ after $2$ steps                                                    \\
            Lukasiewicz $T^s_P$                                                                                                                                                                  & $a^{.8}\lor b^{.8}$; $c^{.8}\lor b^{.8}$, $c^{.8}\lor a^{.8}$; $c^{.8}$                                                                                       & $c^{.8}\in\mathit{MS}_P$                                                            \\
            Guller $\mathcal{C}_P$ \cite{Guller2004procedural}                                                                                                                                   & $(\{a,b\},.8)$; $(\{c,b\},.8)$, $(\{c,a\},.8)$; $(\{c\},.8)$                                                                                                  & $(\{c\},.8)\in\mathit{FS}(P)$                                                       \\
            Nieves $T$ \cite{NievesOC13}                                                                                                                                                         & $P$; $+0.8:c\lor b$, $+0.8:c\lor a$; $+0.8:c$; $P_{\{a,c\}}$ gives $\{(a,.8),(c,.8)\}$                                                                        & $\mathit{Sem}_{\min}\ni(c,.8)$; two answer sets                                     \\
            Grid translation \cite{MushthofaSD14}                                                                                                                                                & $k=5$: $a_i\lor b_i\leftarrow\top$ for $i\leq4$, $c_i\leftarrow a_i$, $c_i\leftarrow b_i$                                                                     & $M_5(A)$: both minimal models                                                       \\
            With $\flpor_l$ instead of $\flpor_g$                                                                                                                                                & $a+b\geq0.8$, $c\geq\max(a,b)$; shift $a\flpif{}0.8\flpand_l\flpnot b$, $b\flpif{}0.8\flpand_l\flpnot a$ (HCF)                                                & continuum $(x,0.8-x,\max(x,0.8-x))$                                                 \\
            Loyer--Straccia \cite{LoyerS02}                                                                                                                                                      & $f_d$ merges the rules of one head atom, not the atoms of one head                                                                                            & not applicable                                                                      \\
        \end{longtable}
    }

\subsection*{Literature examples in our framework}
We ran the examples of the cited papers through the non-deterministic approximator $N$ on grids $\{0,\frac1k,\dots,1\}$ (all stable pairs and the Kripke-Kleene and well-founded state by enumeration).

    {\scriptsize
        \begin{longtable}{@{}p{0.15\textwidth}p{0.27\textwidth}p{0.19\textwidth}p{0.24\textwidth}p{0.05\textwidth}@{}}
            \caption{Examples of related work under $N$.}\label{tab:examples-literature}                                                                                                                                                                                                                                                                                                                                                                                   \\
            \toprule
            \textbf{Source}                                      & \textbf{Program}                                                                                                                                                  & \textbf{Published}                                                    & \textbf{Stable fixpoints $(I,I)$ of $N$}                                                                                        & \textbf{Match}            \\
            \midrule\endfirsthead
            \toprule
            \textbf{Source}                                      & \textbf{Program}                                                                                                                                                  & \textbf{Published}                                                    & \textbf{Stable fixpoints $(I,I)$ of $N$}                                                                                        & \textbf{Match}            \\
            \midrule\endhead
            \bottomrule\endfoot
            \cite{MushthofaSD15}, Ex.~1                          & $\{a\leftarrow\flpnot  c,\ b\leftarrow\flpnot  c,\ c\leftarrow a\flpor_l b,\ d\flpor_l e\leftarrow c\}$                                                           & $(\frac13,\frac13,\frac23,\frac23-x,x)$, none on $L_k$ for $3\nmid k$ & $k=3$: $x\in\{0,\frac13,\frac23\}$; $k=2$: none                                                                                 & yes                       \\
            \cite{MushthofaSD15,MushthofaSD20}, Ex.~2            & $\{a\flpor_l b\leftarrow1,\ a\leftarrow b,\ b\leftarrow a\}$                                                                                                      & $(\frac12,\frac12)$ on $[0,1]$, $(1,1)$ on $L_1$                      & $k=1$: $(1,1)$; $k=2,4$: $(\frac12,\frac12)$, also the WF state                                                                 & yes                       \\
            \cite{MushthofaSD20}, Ex.~3                          & $\{a\flpor_l b\leftarrow1,\ b\leftarrow a,\ a\leftarrow a\flpor_l b\}$                                                                                            & $(1,1)$; its shift has none                                           & $(1,1)$ for $k=1,2,4$; shift: none                                                                                              & yes                       \\
            \cite{BlondeelSVD14}, Ex.~2                          & $\{a\leftarrow\flpnot  b,\ b\leftarrow\flpnot  a\}$                                                                                                               & $(x,1-x)$                                                             & all grid points $(\frac ik,1-\frac ik)$; KK $=$ WF $=((0,0),(1,1))$                                                             & yes                       \\
            \cite{SchockaertJV12}, Ex.~9                         & $\{a\flpor_l b\flpor_l b\leftarrow1\}$ vs.\ with $\flpor_g$                                                                                                       & fuzzy and crisp models differ                                         & $\flpor_l$ ($k=4$): $(0,\frac12),(\frac12,\frac14),(1,0)$; $\flpor_g$: $(1,0),(0,1)$                                            & yes                       \\
            \cite{NievesOC13}, Ex.~5                             & $\{0.6{:}a\lor b,\ 0.4{:}a\leftarrow\flpnot  b,\ 0.8{:}b\leftarrow\flpnot  a\}$, Gödel $\flpnot $                                                                 & $\{(a,.6)\}$, $\{(b,.8)\}$                                            & $(.6,0)$, $(0,.8)$                                                                                                              & yes                       \\
            \cite{NievesOC13}, Ex.~7 (weights on the $0.2$ grid) & $\{.8{:}a\lor b\leftarrow\flpnot  c,\ .6{:}c\leftarrow\flpnot  a,\flpnot  b,\ .8{:}a\leftarrow b,\ 1{:}e\leftarrow b,\ .6{:}b\leftarrow a,\ .4{:}b\leftarrow a\}$ & $\{(a,.7),(b,.6),(e,.6)\}$, $\{(c,.6)\}$ (unscaled)                   & $(a,b,c,e)=(.8,.6,0,.6)$, $(0,0,.6,0)$                                                                                          & yes                       \\
            \cite{BautersSDV14tplp}, Ex.~7                       & $\{.8{:}a;b,\ .6{:}c\leftarrow a,\ .4{:}c\leftarrow b\}$                                                                                                          & $\{a^{.8},c^{.6}\}$, $\{b^{.8},c^{.4}\}$                              & $(.8,0,.6)$, $(0,.8,.4)$                                                                                                        & yes                       \\
            \cite{BautersSDV14tplp}                              & $\{1{:}a\leftarrow\flpnot  a\}$                                                                                                                                   & $\{a^{.5}\}$                                                          & $a=\frac12$                                                                                                                     & yes                       \\
            \cite{HeyninckAB22arxiv}, Ex.~21                     & $\{p\lor q\lor r\leftarrow;\ p\leftarrow\flpnot  q;\ r\leftarrow\flpnot  p;\ q\leftarrow\flpnot  r\}$ (crisp)                                                     & no stable fixpoint, no KK interpretation                              & no stable pair; KK state exists with minimal lower elements $\{p\},\{q\},\{r\}$                                                 & yes                       \\
            Approximator section, Ex.~2 ($P_1$, $P_2$)           & $\{a\flpor_l b\leftarrow1\}$ vs.\ $\{a\flpor_g b\leftarrow1\}$                                                                                                    & ---                                                                   & $a+b=1$ (all grid points) vs.\ $(1,0),(0,1)$                                                                                    & ---                       \\
            Approximator section, Ex.~3                          & $\{a\lor b,\ a\lor c,\ b\lor c\leftarrow1\}$                                                                                                                      & three crisp minimal models                                            & crisp: the three; $\flpor_l$ ($k=4$): also $(\frac12,\frac12,\frac12)$ and $(\frac ik,1-\frac ik,\dots)$; $\flpor_g$: the three & ---                       \\
            \cite{NievesOC13}, counterexample                    & $\{0.9{:}a\lor b,\ 0.5{:}a\}$                                                                                                                                     & only $\{(a,.9)\}$                                                     & $(.9,0)$ and $(.5,.9)$                                                                                                          & no: more stable fixpoints \\
        \end{longtable}

        \begin{longtable}{@{}p{0.15\textwidth}p{0.30\textwidth}p{0.27\textwidth}p{0.24\textwidth}@{}}
            \caption{Where the syntax, the connectives and the semantics are defined.}\label{tab:syntax-locations}                                                                                                     \\
            \toprule
            \textbf{Approach}           & \textbf{Syntax}                                       & \textbf{Connectives}                            & \textbf{Semantics}                                                 \\
            \midrule\endfirsthead
            \toprule
            \textbf{Approach}           & \textbf{Syntax}                                       & \textbf{Connectives}                            & \textbf{Semantics}                                                 \\
            \midrule\endhead
            \bottomrule\endfoot
            \cite{VanNieuwenborghDV06}  & Sect.~2.2, p.~4; Sect.~3, p.~5; Defs.~1--3, pp.~6--7  & Sect.~2.1, p.~3; Ex.~2, p.~4                    & Def.~4, p.~7; Def.~5, pp.~8--9                                     \\
            \cite{VanNieuwenborghDV07}  & Sect.~2.2, p.~368; Sect.~3, pp.~369--370              & Sect.~2.1, pp.~366--367                         & Def.~4, p.~374; Sect.~5                                            \\
            \cite{JanssenSVD09}         & Sect.~2, pp.~353--354                                 & Sect.~2, p.~353; footnote~1, p.~354             & Sect.~3, p.~356; Def.~1, p.~357                                    \\
            \cite{JanssenVSD12}         & Sect.~3, pp.~6--7                                     & Sect.~2, Tables~1--2, p.~5                      & Defs.~1--3, pp.~7--8                                               \\
            \cite{JanssenSVD12core}     & Sect.~3.1, Defs.~1--3, p.~4; Sect.~4.1, eq.~(3), p.~6 & Sect.~2, Tables~1--2, pp.~3--4; Def.~2, p.~4    & Defs.~4--7                                                         \\
            \cite{SchockaertJVD09}      & Sect.~2.1, p.~138; Sect.~2.2, pp.~139--140            & Sect.~2.2, pp.~139--140; Sect.~3.1, p.~141      & Defs.~6--8, p.~142; Props.~2--4                                    \\
            \cite{SchockaertJV12}       & Sect.~2.4, eq.~(10), pp.~120--121; Sect.~3.1          & Sect.~2.3, Table~I, pp.~118--120; p.~121        & Sect.~2.4, p.~121; Defs.~3--4, Sect.~3.1                           \\
            \cite{BlondeelSVD14}        & Sect.~2.1, p.~8; Sect.~2.3, pp.~11--12                & Sect.~2.2, pp.~9--10; Sect.~2.3, pp.~11--12     & Sect.~2.3, pp.~12--13                                              \\
            \cite{MushthofaSD14}        & Sect.~3, PDF p.~2 (right column)                      & Sect.~3, PDF p.~2; Sect.~4.1                    & Sect.~3, PDF p.~2                                                  \\
            \cite{MushthofaSD15}        & Sect.~2.1, pp.~455--456                               & Sect.~2.1, p.~456                               & Sect.~2.1, p.~456                                                  \\
            \cite{MushthofaSD20}        & Sect.~1.3, p.~4                                       & Sect.~1.3, p.~5                                 & Sect.~1.3, p.~5                                                    \\
            \cite{AlvianoP13}           & Sect.~2, p.~3                                         & Sect.~2, p.~3                                   & Sect.~2, p.~3; Def.~1                                              \\
            \cite{AlvianoP15}           & Sect.~2.1, pp.~3--4                                   & Sect.~2.1, p.~4; Sect.~3, pp.~5--6              & Sect.~2.1, p.~4                                                    \\
            \cite{AlvianoAP17}          & Preliminaries, p.~3695                                & Preliminaries, p.~3695                          & Preliminaries, p.~3695                                             \\
            \cite{Lukasiewicz01}        & Sect.~2.3, p.~339                                     & Sect.~2.2, eq.~(1), p.~339                      & Sect.~2.3, p.~340; Sects.~3--4                                     \\
            \cite{LukasiewiczS07}       & Sect.~4.1, eq.~(7), p.~7                              & Sect.~2, p.~2; Tables~1--3                      & Sect.~4.2, eq.~(8), p.~9                                           \\
            \cite{LukasiewiczS09}       & Sect.~4.1, eq.~(20), p.~844                           & Sect.~3.1, Table~1, p.~840                      & Sect.~4.2, p.~845                                                  \\
            \cite{Mateis00}             & Sect.~3.1, pp.~227--228                               & Sect.~2, pp.~226--227                           & Sect.~3.2, eq.~(1), p.~228                                         \\
            \cite{Guller02}             & Sect.~3, p.~253                                       & Table~1, p.~253                                 & Sect.~4, Defs.~2--4, p.~254                                        \\
            \cite{Guller2004procedural} & pp.~519--520; Sect.~2.2                               & Sect.~2.1, p.~520; Table~1, p.~521              & Sect.~3, Defs.~1--3, pp.~523--524                                  \\
            \cite{Saad2009extended}     & Sect.~4, Defs.~1--2, pp.~228--229                     & Sect.~3, pp.~227--228                           & Sect.~4.1, Defs.~3--9, pp.~230--231; Sect.~5                       \\
            \cite{Saad2010disjunctive}  & Sect.~3, Defs.~1--2, p.~309                           & Sect.~2, pp.~308--309                           & Sect.~3.1, Defs.~3--9, pp.~310--312; Sect.~4, Defs.~10--11, p.~313 \\
            \cite{NievesOC07}           & Sect.~2.2, Remark~1, p.~8                             & Def.~4, p.~9                                    & Sect.~4, Defs.~5--6, pp.~15--16                                    \\
            \cite{NievesOC13}           & Sect.~2.2, Remark~1, p.~40                            & Def.~4, p.~41                                   & Sect.~4, Defs.~5--6, pp.~47--48                                    \\
            \cite{NicolasGSL06}         & Defs.~2--3                                            & Defs.~8--9, p.~150                              & Defs.~10--11, p.~152                                               \\
            \cite{BautersSDV14tplp}     & Sect.~2.1, p.~6; Sect.~3                              & Sect.~4, eqs.~(5)--(6), pp.~19--20              & Defs.~4--5, p.~15; Sect.~4                                         \\
            \cite{BautersSDV14ijar}     & Sect.~2.1, p.~6                                       & Sect.~2.3, p.~11                                & Defs.~1--2, p.~12                                                  \\
            \cite{StracciaOD09}         & Sect.~4, p.~1898                                      & Sect.~4, p.~1898; Props.~4.1--4.5               & Sect.~4.1, p.~1901                                                 \\
            \cite{LoyerS09}             & Sect.~2, p.~391                                       & Sect.~2, pp.~390--391                           & Defs.~1--3, pp.~393--394                                           \\
            \cite{MadridO12}            & Def.~2.2, p.~311                                      & Def.~2.1, Remark~1, p.~311                      & Defs.~2.3--2.4, p.~312                                             \\
            \cite{MedinaOV01}           & Sect.~3.1, Def.~10, p.~356                            & Defs.~4--6, pp.~352--354                        & Defs.~11--14, pp.~356--357                                         \\
            \cite{CornejoLM18}          & Sect.~3, Def.~17, p.~9                                & Defs.~9--11, pp.~6--7                           & Def.~18, p.~10; Def.~20, p.~11; Def.~24, p.~13                     \\
            \cite{CornejoLM20}          & Sect.~3, Defs.~9--10, pp.~8--9                        & Def.~9, pp.~8--10                               & Defs.~12--14, pp.~11--12                                           \\
            \cite{CornejoLM20math}      & Def.~2, p.~3; Defs.~6--7, p.~5                        & Def.~1, pp.~2--3                                & Defs.~3--5, pp.~3--5; Defs.~8--12                                  \\
            \cite{HeyninckAB22arxiv}    & Sect.~2.1, p.~5                                       & Sect.~2.1, pp.~5--6; Def.~3                     & Def.~1, p.~6; Defs.~2--3                                           \\
            \cite{PelovT04}             & Preliminaries, PDF pp.~1--2                           & Preliminaries, PDF pp.~1--2                     & Defs.~1--5, PDF pp.~2--4                                           \\
            \cite{RiegelEtAl20}         & Sect.~3, pp.~3--4; eqs.~(1)--(3), p.~4                & Sect.~3.1, pp.~3--4; eqs.~(9)--(24), pp.~14--15 & Sect.~4, pp.~5--6                                                  \\
            \cite{SenCRG22}             & Sect.~4, PDF pp.~3--4                                 & Sect.~3, PDF pp.~2--3; eq.~(1)                  & Sects.~4.1--4.4, PDF pp.~4--5                                      \\
            \cite{Antic24}              & Sect.~3.1, rule~(3), p.~5                             & Sect.~2.1, p.~2; Def.~7                         & Defs.~1--5, pp.~3ff.                                               \\
            \cite{KettmannHS25}         & Sect.~2.2, eqs.~(1)--(2), PDF p.~3                    & Sect.~2.2, PDF p.~3                             & Eq.~(4), Def.~1, PDF pp.~3--4                                      \\
            \cite{KettmannSHS26arxiv}   & Sect.~3, Def.~1, pp.~6--7                             & Sect.~2.2, p.~5; p.~7                           & Defs.~3--4, pp.~8 and 10                                           \\
        \end{longtable}
    }
\endgroup
